% =====================================================================================
% Reproduction — the short version. The long one is Appendix C, which this points at.
% No number appears here, so no macro is cited.
% =====================================================================================
\chapter*{Reproducing this book}
\addcontentsline{toc}{chapter}{Reproducing this book}
\markboth{Reproducing this book}{Reproducing this book}

Everything in these pages is runnable. The analyses live in fourteen notebooks; the book is
generated from them. \Cref{app:env} gives the full account --- exact commands, the notebook-to-chapter
map, and what each environment variable changes. This page is the short one.

\section*{Two environments, and why}

The code ships \emph{two} Python environments, and the split is the single largest engineering
decision behind the repository, so it is stated rather than hidden.

\begin{description}[leftmargin=1.2em, style=nextline, font=\normalfont\bfseries]
  \item[\code{cmp-core} --- pymc $\ge$ 6.] The foundations, the uplift and hierarchical models,
  mediation, the choice of controls, and synthetic control --- plus the interactive apps and the
  package tests.
  \item[\code{cmp-legacy} --- pymc $<$ 6.] The chapters built on \code{causalpy} and
  \code{pymc-marketing}: media mix, difference-in-differences, regression discontinuity,
  interrupted time series, and instrumental variables.
\end{description}

The reason is not taste. As of writing, \code{pymc-bart}, \code{bambi} and \code{pathmc} require
pymc~6, while \code{causalpy} and \code{pymc-marketing} still pin pymc~5. No single environment
satisfies both. Holding the whole book back to pymc~5 would have cost the modern BART and
path-analysis work; vendoring patched forks would have made the results unreproducible for the
reader. Two environments is the honest cost of a genuinely split ecosystem, and when the two
libraries move to pymc~6 the split collapses back into one.

\section*{Two switches}

\code{CMP\_FAST} (default \code{1}) runs every notebook in a teaching mode: short chains, few
trees, under two minutes, suitable for a live rerun in front of a class. \textbf{Every number in
this book comes from the full run, \code{CMP\_FAST=0}.} The fast mode is for teaching; it is not
what anything here was graded on.

\code{CMP\_REAL} (default \code{0}) gates the cells that fetch a real public dataset over the
network, so that the offline test suite stays deterministic. It is set for the chapters that grade
themselves against real data: \cref{ch:foundations} and \cref{ch:ctrl} on the LaLonde job-training
benchmark, \cref{ch:geo} on Google's published geo experiment, and \cref{ch:iv} on Criteo's
advertising experiment.

\section*{Rebuilding the book}

\begin{quote}
\code{make book}
\end{quote}

That regenerates the macro file from the notebooks' emitted results, re-renders every figure and
table, and typesets the book and the per-chapter offprints.

One consequence of the injection rule (\emph{A note on the numbers}, above) deserves to be said
plainly, because it will otherwise be met as a bug. \textbf{A fresh clone cannot build the book.}
The generated directory holds no results until the notebooks have been executed, and a chapter that
cites a result the notebooks have not emitted does not typeset a blank --- it halts the build. So
the first step in any reproduction is to run the notebooks at full quality, in both environments,
and only then to typeset. That the book refuses to compile from an empty results directory is the
same property, seen from the other side, that guarantees no number in it is stale.
